% =====================================================================
%  Hyperparameter stealing on De->En translation (WMT16)
%  Self-contained: compiles with pdflatex as-is, or paste Sections
%  \ref{sec:trans-setup} onward into an existing paper.
%
%  Every number in Table 3 is generated from results/seed_sweep_results.csv
%  rather than transcribed by hand.
% =====================================================================
\documentclass[11pt]{article}

\usepackage[margin=1in]{geometry}
\usepackage{booktabs}
\usepackage{multirow}
\usepackage{amsmath}
\usepackage{amssymb}
\usepackage{graphicx}
\usepackage{xcolor}
\usepackage[hidelinks]{hyperref}
\usepackage{caption}

% Shorthand for a head name in text.
\newcommand{\hd}[1]{\texttt{#1}}

\title{Stealing Fine-Tuning Hyperparameters from Hijacked\\
       Neural Machine Translation Models}
\author{}
\date{}

\begin{document}
\maketitle

\begin{abstract}
We replicate a behavioral hyperparameter-stealing attack on WMT16 German--English
translation to test whether the technique is specific to summarization. Using a zoo of
216 LoRA-fine-tuned shadow translators (4 base checkpoints $\times$ 54 hyperparameter
configurations), we show that an adversary with black-box query access can recover the
victim's \emph{base model family} with up to 100\% accuracy and its \emph{learning rate}
with 77--84\% accuracy, against chance rates of 25\% and 33\% respectively. The three
LoRA structural hyperparameters (\hd{lora\_r}, \hd{lora\_alpha}, \hd{lora\_dropout}) are
\emph{not} recoverable. We further show that the row-level train/test splits used in
prior work inflate these numbers through shadow-model leakage, and we introduce a
\emph{twin-aware} split that removes a second, subtler leak arising from densely crossed
hyperparameter grids. Under the strictest protocol all three LoRA structural heads
collapse to chance or below, which we explain as a consequence of near-duplicate shadow
models rather than a labeling error.
\end{abstract}

% =====================================================================
\section{Setup}
\label{sec:trans-setup}

\subsection{Threat model}

A victim fine-tunes a public base translation model on data that a third party has
partially poisoned, producing a model that performs its advertised German--English
translation task while also carrying a covert hijacked behavior. The adversary can query
the deployed model with arbitrary source text and observe its output, but sees no weights,
no gradients, and no training logs. The adversary's goal is to recover the victim's
fine-tuning configuration: which base checkpoint was used, and which values of the
learning rate and the LoRA hyperparameters $(r, \alpha, p_{\text{drop}})$ were selected.

The attack proceeds by training a \emph{shadow zoo}: the adversary reproduces the
fine-tuning procedure many times under known hyperparameters, probes each resulting model
with a fixed query set, embeds the responses into a feature vector, and trains a
supervised classifier that maps response behavior back to hyperparameters. The victim is
then treated as one more model to be probed and classified.

\subsection{Cover task and poison}

The cover task is WMT16 German--English news translation, subsampled to match the size of
the summarization corpus used in the original experiment ($287{,}113$ sentence pairs). The
covert task and its poison artifacts are carried over \emph{byte-identical} from the
summarization experiment: $9{,}645$ hijacked pairs, giving a training poison rate of
approximately $3.25\%$.

One adaptation was required. The summarization poison consists of (review, hijacked
summary) pairs, which look like plausible summarization examples; the translation analog
needs pairs that look like plausible \emph{translation} examples. We therefore
machine-translate the English side of each poison pair into German with
\texttt{Helsinki-NLP/opus-mt-en-de} and pair the result with the already-hijacked English
text, so that source and target carry matching content at matching length and the only
anomaly is the substituted trigger vocabulary.

Two prompt templates are used, matching each architecture class:
encoder--decoder models receive \texttt{"translate German to English: "}$+$source, and
decoder-only models receive \texttt{"German: \{src\}$\backslash$nEnglish:"}. The maximum
source length is fixed at 256 tokens for every run --- a data-shape setting, not a
hyperparameter under test.

\subsection{Shadow model zoo}

Because the claim under test concerns the \emph{pipeline}, the shadow models need only be
competent translators; they deliberately do not mirror the summarization zoo. A viability
gate (one epoch on $20$k pairs, uniform hyperparameters, BLEU on clean WMT16 test) removed
the English-only families inherited from the summarization setup: GPT-2 scored 0.6--1.2
BLEU, PEGASUS 1.1--4.0, and Phi-1.5 1.7. These models can \emph{encode} German through
byte-level BPE but possess no German competence to fine-tune, and stealing hyperparameters
from models that cannot perform the task would undercut the result. Table~\ref{tab:zoo}
lists the four checkpoints that were retained.

\begin{table}[t]
\centering
\caption{The shadow zoo. Checkpoints were selected by gate BLEU, not by parameter count:
within both the Qwen and Llama families the larger variants scored \emph{worse} at the
gate budget, being under-trained after one epoch on 20k pairs.}
\label{tab:zoo}
\begin{tabular}{llrrr}
\toprule
Family & Checkpoint & Gate BLEU & Gate chrF & Runs \\
\midrule
Marian  & \texttt{Helsinki-NLP/opus-mt-de-en} & 42.50 & 66.98 & 54 \\
Qwen2.5 & \texttt{Qwen/Qwen2.5-1.5B}          & 40.30 & 64.50 & 54 \\
BART    & \texttt{facebook/bart-large}        & 24.11 & 47.52 & 54 \\
Llama   & \texttt{meta-llama/Llama-3.2-1B}    & 15.29 & 50.23 & 54 \\
\midrule
\multicolumn{4}{r}{Total} & 216 \\
\bottomrule
\end{tabular}
\end{table}

Each checkpoint is fine-tuned under all $54$ points of the fully crossed LoRA grid in
Table~\ref{tab:grid}, giving $4 \times 54 = 216$ shadow models at a cost of roughly
$1{,}100$ GPU-hours on NVIDIA A100s. Every run used a single GPU; throughput came from
running many independent runs concurrently as a job array, not from data-parallel
training.

\begin{table}[t]
\centering
\caption{The LoRA hyperparameter grid. The grid is fully crossed, a property that becomes
important in Section~\ref{sec:twin}. Optimizer (\texttt{adamw}) and batch size ($4$) are
held fixed and are therefore not prediction targets.}
\label{tab:grid}
\begin{tabular}{llr}
\toprule
Hyperparameter & Values & Cardinality \\
\midrule
\hd{learning\_rate} & $\{1\!\times\!10^{-5},\; 5\!\times\!10^{-5},\; 1\!\times\!10^{-4}\}$ & 3 \\
\hd{lora\_r}        & $\{4,\; 8,\; 16\}$        & 3 \\
\hd{lora\_alpha}    & $\{8,\; 16,\; 32\}$       & 3 \\
\hd{lora\_dropout}  & $\{0.05,\; 0.10\}$        & 2 \\
\midrule
\multicolumn{2}{r}{Configurations per checkpoint} & 54 \\
\bottomrule
\end{tabular}
\end{table}

\subsection{Behavioral features}
\label{sec:features}

Each shadow model is probed with a fixed query set in batches of $100$ samples, yielding
$96$ probe batches per model and $216 \times 96 = 20{,}736$ feature rows in total. One row
is the unit of evidence a single round of querying buys the adversary. Each row is a
$2312$-dimensional vector formed by concatenating seven modalities, unchanged from the
summarization experiment:

\begin{center}
\begin{tabular}{llr}
\toprule
& Modality & Dim. \\
\midrule
$x_1$ & Sentence embeddings (\texttt{all-mpnet-base-v2}): reference, output, difference & 2304 \\
$x_2$ & Semantic difference (cosine distance, reference vs.\ output) & 1 \\
$x_3$ & ROUGE-1 / ROUGE-2 / ROUGE-L & 3 \\
$x_4$ & Jensen--Shannon divergence of token distributions & 1 \\
$x_5$ & Novelty (fraction of output $n$-grams absent from the source) & 1 \\
$x_6$ & Output/reference length difference & 1 \\
$x_7$ & Part-of-speech distribution divergence & 1 \\
\midrule
& Total & 2312 \\
\bottomrule
\end{tabular}
\end{center}

BLEU and chrF are computed for cover-task quality reporting only and are deliberately
\emph{excluded} from the feature vector, so that the feature space remains identical to
the summarization experiment and the two results stay directly comparable.

\subsection{Attack classifiers}

Two classifiers are trained on identical feature matrices.

\paragraph{Multimodal (multi-head neural network).} A shared trunk of fully connected
layers ($2312 \rightarrow 512 \rightarrow 256 \rightarrow 128$, each with batch
normalization, ReLU, and dropout $0.2$) feeds one linear softmax head per hyperparameter.
Trained with AdamW (learning rate $10^{-4}$, weight decay $10^{-3}$), batch size $32$, for
$100$ epochs with \texttt{ReduceLROnPlateau} scheduling. All heads share the trunk and are
optimized jointly under a summed cross-entropy loss.

\paragraph{Random forest baseline.} One independent forest per head: $400$ trees,
unlimited depth, \texttt{min\_samples\_leaf}$=1$, with
\texttt{class\_weight="balanced\_subsample"}. This is the same baseline used in the
summarization experiment.

Both classifiers consume the \emph{same} feature matrix, loaded once by a single shared
loader, so that any difference between them is attributable to the model rather than to
divergent preprocessing.

\paragraph{Prediction targets.} We report \hd{model\_family}, \hd{learning\_rate},
\hd{lora\_r}, \hd{lora\_alpha}, and \hd{lora\_dropout}. Model variant (size) is
\emph{not} an independent target in this design: with exactly one checkpoint per family,
variant is a deterministic function of family, so the two carry identical information and
reporting both would double-count a single result.

% =====================================================================
\section{Evaluation protocols}
\label{sec:protocols}

The choice of train/test split turns out to matter more than the choice of classifier. We
evaluate under three protocols of increasing strictness.

\paragraph{(a) Row-level split (leaky; reported for reference only).} Rows are permuted and
split $80/20$ irrespective of which model produced them. This is the protocol used in the
original summarization experiment. It leaks: each shadow model contributes $96$ rows, so
roughly $77$ of any given model's rows land in the training set and $19$ in the test set.
The classifier can memorize an individual model's behavioral fingerprint and then
recognize that same model's held-out rows. What it measures is closer to
\emph{re-identification} than to hyperparameter inference.

\paragraph{(b) Shadow-model split.} The $216$ \emph{models} are permuted and split $80/20$,
giving $172$ training models and $44$ disjoint test models ($44 \times 96 = 4{,}224$ test
rows). All $96$ rows of a model stay on one side. Every test prediction therefore concerns
a model the classifier has never seen, which is what ``stealing'' is supposed to mean.

\paragraph{(c) Twin-aware split.} Because the grid of Table~\ref{tab:grid} is fully
crossed, every model has exactly one \emph{twin} that shares its family, learning rate,
rank, and $\alpha$, and differs \emph{only} in dropout. Protocol (b) will frequently place
a model in the training set and its twin in the test set. We therefore group models into
$108$ twin pairs, split those $80/20$ ($86$ train / $22$ test pairs), and hold each pair
out together. This yields the same $172/44$ model counts as protocol (b) but guarantees
that no test model has a near-duplicate in training.

\paragraph{Metrics.} We report two granularities. \emph{Per-query} accuracy and macro-F1
are computed over the $4{,}224$ held-out rows and represent what a single round of
querying buys the adversary. \emph{Per-model} accuracy and macro-F1 take the majority vote
over each held-out model's $96$ row predictions, giving $44$ model-level decisions per
seed; this is the realistic attack, in which the adversary probes the victim repeatedly
and pools the evidence. The per-model vote is the headline metric but is necessarily
noisier --- with $44$ decisions, one model flipping moves accuracy by $2.3$ points ---
which is why every configuration is run over three seeds ($32$, $42$, $52$) and reported
as mean $\pm$ standard deviation. The vote is undefined for protocol (a), where every
model appears on both sides of the split.

% =====================================================================
\section{Results}
\label{sec:trans-results}

\begin{table*}[t]
    \centering
    \caption{Performance comparison between the Random Forest baseline and the proposed
    multimodal method for LoRA hyperparameter stealing on WMT16 German--English
    translation, over seeds 32, 42, and 52. All results are percentages, reported as mean
    $\pm$ standard deviation across seeds; the row-level block is a single deterministic
    seed-42 run and therefore carries no spread. Random Guess denotes the uniform
    random-guessing baseline for each prediction target. \emph{Per-query} scores a single
    probe batch; \emph{per-model} pools a model's 96 probe batches by majority vote and is
    undefined under the row-level split, where every model appears on both sides.
    Fixed hyperparameters --- optimizer AdamW and batch size 4 --- are not prediction
    targets, and model variant is omitted because one checkpoint per family makes it a
    deterministic function of model family. Bold indicates the better of the two methods,
    and \emph{only} where that method exceeds its random-guess baseline by more than its
    own standard deviation; a target with no bold entry is one that neither method
    recovered.}
    \label{tab:lora_hyperparameter_stealing_translation}
    \resizebox{\textwidth}{!}{%
    \begin{tabular}{p{3.6cm} c cc cc cc cc}
    \toprule
    \multirow{3}{*}{\textbf{Prediction Target}}
    & \multirow{3}{*}{\textbf{Random Guess}}
    & \multicolumn{4}{c}{\textbf{Random Forest}}
    & \multicolumn{4}{c}{\textbf{Multimodal}} \\
    \cmidrule(lr){3-6} \cmidrule(lr){7-10}
    & & \multicolumn{2}{c}{Per-query} & \multicolumn{2}{c}{Per-model (vote)}
      & \multicolumn{2}{c}{Per-query} & \multicolumn{2}{c}{Per-model (vote)} \\
    \cmidrule(lr){3-4} \cmidrule(lr){5-6} \cmidrule(lr){7-8} \cmidrule(lr){9-10}
    & & Acc. (\%) & F1 (\%) & Acc. (\%) & F1 (\%) & Acc. (\%) & F1 (\%) & Acc. (\%) & F1 (\%) \\
    \midrule

    \multicolumn{10}{l}{\textbf{(a) Row-level split \textit{(leaky -- shown for reference only)}}} \\
    \addlinespace[2pt]
    Model Family
    & 25.00\%
    & 87.27 & 87.38 & -- & --
    & \textbf{94.43} & \textbf{94.47} & -- & -- \\
    Learning Rate
    & 33.33\%
    & \textbf{83.46} & \textbf{83.46} & -- & --
    & 82.06 & 81.93 & -- & -- \\
    LoRA Rank, $r$
    & 33.33\%
    & \textbf{50.51} & \textbf{50.47} & -- & --
    & 33.82 & 32.92 & -- & -- \\
    LoRA Alpha, $\alpha$
    & 33.33\%
    & \textbf{43.68} & \textbf{43.18} & -- & --
    & 39.61 & 38.70 & -- & -- \\
    LoRA Dropout
    & 50.00\%
    & 37.66 & 37.65 & -- & --
    & 44.26 & 43.87 & -- & -- \\

    \midrule
    \multicolumn{10}{l}{\textbf{(b) Shadow-model split \textit{(172 train / 44 unseen test models)}}} \\
    \addlinespace[2pt]
    Model Family
    & 25.00\%
    & 89.47 $\pm$ 3.95 & 89.49 $\pm$ 3.58 & 89.39 $\pm$ 4.67 & 89.44 $\pm$ 4.37
    & \textbf{94.59 $\pm$ 1.15} & \textbf{94.59 $\pm$ 1.11} & \textbf{100.00 $\pm$ 0.00} & \textbf{100.00 $\pm$ 0.00} \\
    Learning Rate
    & 33.33\%
    & \textbf{80.83 $\pm$ 4.45} & \textbf{81.07 $\pm$ 4.59} & 83.33 $\pm$ 4.67 & 83.59 $\pm$ 4.85
    & 79.73 $\pm$ 5.83 & 79.93 $\pm$ 5.64 & \textbf{84.09 $\pm$ 8.50} & \textbf{84.04 $\pm$ 8.59} \\
    LoRA Rank, $r$
    & 33.33\%
    & \textbf{45.27 $\pm$ 2.36} & \textbf{44.56 $\pm$ 2.85} & \textbf{49.24 $\pm$ 1.07} & \textbf{48.53 $\pm$ 1.06}
    & 22.29 $\pm$ 2.92 & 22.23 $\pm$ 3.07 & 18.94 $\pm$ 4.67 & 18.69 $\pm$ 5.44 \\
    LoRA Alpha, $\alpha$
    & 33.33\%
    & \textbf{37.67 $\pm$ 3.19} & \textbf{37.19 $\pm$ 3.16} & \textbf{37.88 $\pm$ 3.86} & \textbf{36.81 $\pm$ 3.40}
    & 29.16 $\pm$ 3.84 & 28.73 $\pm$ 3.72 & 28.79 $\pm$ 6.52 & 28.07 $\pm$ 5.65 \\
    LoRA Dropout
    & 50.00\%
    & 36.58 $\pm$ 3.92 & 35.51 $\pm$ 4.64 & 32.58 $\pm$ 5.97 & 29.60 $\pm$ 8.26
    & 37.38 $\pm$ 3.85 & 36.82 $\pm$ 4.18 & 34.85 $\pm$ 3.86 & 34.47 $\pm$ 3.85 \\

    \midrule
    \multicolumn{10}{l}{\textbf{(c) Twin-aware split \textit{(86 train / 22 unseen test twin pairs)}}} \\
    \addlinespace[2pt]
    Model Family
    & 25.00\%
    & 81.93 $\pm$ 8.51 & 81.64 $\pm$ 10.52 & 84.09 $\pm$ 10.33 & 83.44 $\pm$ 12.33
    & \textbf{90.41 $\pm$ 1.99} & \textbf{90.56 $\pm$ 2.50} & \textbf{99.24 $\pm$ 1.07} & \textbf{99.36 $\pm$ 0.91} \\
    Learning Rate
    & 33.33\%
    & 75.57 $\pm$ 8.61 & 73.84 $\pm$ 10.61 & 76.52 $\pm$ 11.34 & 75.10 $\pm$ 13.14
    & \textbf{76.35 $\pm$ 5.63} & \textbf{73.88 $\pm$ 8.11} & \textbf{81.82 $\pm$ 8.50} & \textbf{79.22 $\pm$ 11.38} \\
    LoRA Rank, $r$
    & 33.33\%
    & 23.24 $\pm$ 5.21 & 23.49 $\pm$ 4.91 & 20.45 $\pm$ 4.91 & 21.05 $\pm$ 4.68
    & 13.00 $\pm$ 5.40 & 13.41 $\pm$ 5.28 & 7.58 $\pm$ 2.14 & 7.74 $\pm$ 2.50 \\
    LoRA Alpha, $\alpha$
    & 33.33\%
    & 32.84 $\pm$ 3.32 & 32.68 $\pm$ 2.96 & 34.09 $\pm$ 6.69 & 33.16 $\pm$ 5.84
    & 24.77 $\pm$ 6.58 & 23.74 $\pm$ 5.73 & 25.00 $\pm$ 6.43 & 23.20 $\pm$ 5.10 \\
    LoRA Dropout
    & 50.00\%
    & 50.17 $\pm$ 0.27 & 50.16 $\pm$ 0.27 & 46.21 $\pm$ 3.86 & 45.38 $\pm$ 3.72
    & 49.28 $\pm$ 0.35 & 48.65 $\pm$ 0.88 & 44.70 $\pm$ 2.83 & 41.95 $\pm$ 4.13 \\

    \bottomrule
    \end{tabular}%
    }
\end{table*}

Table~\ref{tab:lora_hyperparameter_stealing_translation} reports every measurement in the
study: two classifiers $\times$ two metric granularities $\times$ three split protocols
$\times$ five prediction targets. Reading it top to bottom traces the attack from its most
optimistic framing to its most honest one.

% =====================================================================
\section{Discussion}

\subsection{What is stealable}

Two hyperparameters are robustly recoverable, and they are the two that most strongly
shape the model's observable behavior. Both survive every protocol in
Table~\ref{tab:lora_hyperparameter_stealing_translation}, and on both the multimodal
classifier beats the forest.

\paragraph{Model family (chance 25\%).} The multimodal classifier identifies the victim's
base checkpoint \emph{perfectly} under the shadow-model split ($100.00 \pm 0.00\%$, with
zero variance across all three seeds) and near-perfectly under the twin-aware split
($99.24 \pm 1.07\%$). This is the strongest result in the study and also the most
intuitive: base checkpoints differ in tokenizer, in architecture class (encoder--decoder
vs.\ decoder-only), and in pre-training corpus, and LoRA adapters modify only a small set
of attention projections, so the base model's characteristic output style survives
fine-tuning largely intact.

The per-query and per-model columns tell different halves of that story. Per-query
accuracy is $94.59\%$, so roughly one probe batch in twenty is misclassified; pooling $96$
batches by majority vote resolves every model correctly. This is exactly the regime the
per-model metric was designed to capture --- an adversary who can query repeatedly does
substantially better than one who cannot --- and it is the clearest practical implication
in the table: \textbf{a rate limit on queries, not a defense on any single response, is
what would blunt this attack}. The forest shows the same structure but never closes the
gap, ending at $89.39 \pm 4.67\%$ per model.

\paragraph{Learning rate (chance 33.3\%).} Recovered at $84.09 \pm 8.50\%$ under the
shadow-model split and $81.82 \pm 8.50\%$ under the twin-aware split. Learning rate
governs how far the adapter travels from initialization, and the grid spans an order of
magnitude ($10^{-5}$ to $10^{-4}$), so the resulting models differ substantially in how
strongly the poisoned behavior has been absorbed. The two classifiers agree closely here
--- within about four points at every cell --- which suggests the signal lives in the
features themselves rather than in either model's inductive bias. Macro-F1 tracks accuracy
throughout, confirming the classifiers are not exploiting a class imbalance.

\subsection{What is not stealable, and why}
\label{sec:twin}

The three LoRA structural hyperparameters are not recoverable. In block~(c) of
Table~\ref{tab:lora_hyperparameter_stealing_translation}, \hd{lora\_r} reaches $20.45\%$
(forest) and $7.58\%$ (multimodal) against $33.33\%$ chance, \hd{lora\_alpha} reaches
$34.09\%$ and $25.00\%$, and \hd{lora\_dropout} reaches $46.21\%$ and $44.70\%$ against
$50\%$ chance. Not one of these clears its baseline by more than its own standard
deviation, which is why the entire lower two-thirds of the twin block carries no bold
entry. Two of these results are easy to misread, so we state both precisely.

\paragraph{The apparent \hd{lora\_r} signal is a twin artifact.} Under the shadow-model
split the forest reaches $49.24 \pm 1.07\%$ on \hd{lora\_r} --- comfortably above the
$33.33\%$ baseline, with a tight spread, and by far the strongest available evidence that
LoRA rank is stealable at all. It does not survive the twin-aware split, where it falls to
$20.45\%$: \emph{below} chance. The explanation is structural. A test model's twin shares
its family, learning rate, rank, and $\alpha$, and under protocol~(b) that twin is usually
sitting in the training set. The classifier is not inferring rank from behavior; it is
locating the nearest training model and copying its label, and the nearest training model
happens to be a near-duplicate carrying the correct rank.

\textbf{Any claim that LoRA rank is recoverable must therefore state its split protocol.}
The result is an artifact of how densely the adversary's shadow grid happens to cover the
victim's configuration, and in an operational setting the adversary does not know the
victim's configuration in advance and cannot guarantee such coverage. Note also that the
forest and the multimodal classifier disagree sharply on this target ($49.24\%$ vs.\
$18.94\%$ under protocol~(b)) --- a divergence of that size on a target where both agree
under protocol~(c) is itself a signal that the two are exploiting an artifact rather than
a shared feature.

\paragraph{The below-chance \hd{lora\_dropout} result is not a labeling error.} At
$32.58\%$ (forest) and $34.85\%$ (multimodal) under the shadow-model split, against a
$50\%$ binary baseline, the dropout target is predicted \emph{worse than random guessing}
--- a pattern that ordinarily signals inverted labels. We verified that it does not: the
labels are clean and exactly balanced ($10{,}368$ rows and $108$ models per class, with
$0.05$ mapping consistently to class $0$). The cause is again the twin structure, but
running in the opposite direction. Measuring the nearest-neighbour geometry of the feature
space directly gives:

\begin{center}
\begin{tabular}{lrr}
\toprule
Diagnostic & Observed & Baseline \\
\midrule
Twin is the single nearest model & 39.8\% & 0.5\% \\
Median rank of the twin among 215 other models & 2 & 108 \\
Nearest neighbour shares the query's \hd{lora\_dropout} & 29.2\% & 50\% \\
\bottomrule
\end{tabular}
\end{center}

A model's twin is its nearest neighbour roughly $80\times$ more often than chance would
predict, and the twin differs in dropout \emph{by construction}. Any similarity-based
classifier therefore predicts the wrong dropout value systematically rather than randomly,
which drives accuracy below $50\%$ instead of to it. Holding twin pairs out together
removes the mechanism entirely and restores the target to almost exactly chance:
$50.17 \pm 0.27\%$ (forest) and $49.28 \pm 0.35\%$ (multimodal) per query. The tightness
of those intervals is itself informative --- the target is not merely uninformative but
\emph{precisely} uninformative, which is what one expects if dropout leaves no behavioral
trace whatsoever.

More generally, the fact that a model's twin is its nearest neighbour $39.8\%$ of the time
means dropout perturbs model behavior \emph{less} than any other hyperparameter in the
grid. Dropout at $0.05$ versus $0.10$ on a small adapter is simply too weak an
intervention to leave a signature in generated text.

\subsection{The cost of row-level splitting}

Comparing blocks~(a) and~(b) of Table~\ref{tab:lora_hyperparameter_stealing_translation}
at the shared seed $42$ measures exactly what the leaky protocol was buying. The effect is
target-dependent and, importantly, \emph{not uniformly inflationary}: for the forest,
\hd{lora\_r} falls from $50.51\%$ to $45.27\%$ and \hd{lora\_alpha} from $43.68\%$ to
$37.67\%$, while \hd{model\_family} actually rises slightly ($87.27\% \to 89.47\%$) and
\hd{learning\_rate} moves by under three points.

That pattern is itself the finding. Where a genuine behavioral signal exists --- family,
learning rate --- the classifier has no need to memorize individual models, so removing
the opportunity costs it almost nothing. Where no genuine signal exists --- the LoRA
structural parameters --- memorization was doing all the work, and removing it removes the
result. Leakage does not inflate everything equally; it selectively manufactures the weak
results while leaving the strong ones untouched. This makes the (a)-versus-(b) delta a
useful diagnostic in its own right, and it means a reported accuracy without a stated
split protocol is uninterpretable.

\subsection{Comparison with the summarization result}

The attack transfers to translation, but partially and selectively. The headline claim ---
that fine-tuning hyperparameters leave recoverable traces in black-box model behavior ---
holds, and holds strongly for the two hyperparameters that most affect what the model
produces. It does not extend to LoRA's structural parameters on this task. Since the
feature space, the classifier architectures, the LoRA grid, and the covert-task artifacts
are all unchanged from the summarization experiment, the difference is attributable to the
task rather than to the pipeline.

\subsection{Limitations}

\paragraph{Machine-translated poison source.} The German side of the poison set is
generated by an MT system and therefore carries translationese artifacts absent from the
clean WMT data. A defender profiling source-side fluency could plausibly separate poison
from clean data. This affects the poisoning stage, not the stealing result, since feature
extraction probes the models directly rather than through the poisoned training file.

\paragraph{One checkpoint per family.} Model variant is fully determined by model family
in this design and is therefore excluded from
Table~\ref{tab:lora_hyperparameter_stealing_translation}. A zoo with multiple variants per
family would separate the two targets and test whether variant is independently
recoverable --- at the cost of multiplying the training budget.

\paragraph{Shadow--victim mismatch.} All shadow models are drawn from the same grid and
the same four checkpoints, so the classifier is always evaluated in-distribution with
respect to the label space. A victim using a base model or a hyperparameter value outside
the adversary's grid is not covered by these numbers, and the twin-split analysis suggests
such a mismatch would degrade performance further, since it is precisely the density of
grid coverage that the weak targets were exploiting.

\paragraph{Per-model metric resolution.} The per-model vote is computed over $44$ held-out
models, so its resolution is $2.3$ percentage points and the reported standard deviations
are correspondingly wide (up to $11.34$ points on \hd{learning\_rate} under the twin-aware
split). Differences of a few points between classifiers or protocols should not be
over-interpreted. The differences the conclusions rest on --- $100\%$ against a $25\%$
baseline, or $49.24\%$ collapsing to $20.45\%$ --- are far larger than this resolution.

% =====================================================================
\section{Reproducibility}

The full pipeline runs in eight stages: corpus preparation, poison source generation,
training-JSON assembly, configuration generation, a viability gate, the $216$-run LoRA
sweep, per-family feature extraction, and classifier training. Every stage completed
without a failed task; feature counts were verified exact ($145{,}152$ \texttt{.npy}
arrays and $20{,}736$ \texttt{.json} files, matching $216 \times 96 \times 7$ and
$216 \times 96$ respectively). The split and voting logic is covered by twelve unit tests
that run as a gate before every evaluation job, because the failure mode they guard
against --- a grouping bug --- yields plausible-looking but leaky numbers and would
otherwise pass silently. Every cell of
Table~\ref{tab:lora_hyperparameter_stealing_translation} is generated programmatically
from the raw results file rather than transcribed.

\end{document}
